\chapter{Work Plan and Conclusion}
\label{ch:plan}

\section{Work Plan}
\label{sec:plan-work}

The work takes four months (Table~\ref{tab:plan}). Everything before the
registration point prepares the experiments; everything after it runs them as
planned.

\begin{table}[htbp]
\centering
\caption{Work plan.}
\label{tab:plan}
\footnotesize
\begin{tabularx}{\textwidth}{@{}>{\raggedright\arraybackslash}p{0.1\textwidth} >{\raggedright\arraybackslash}X >{\raggedright\arraybackslash}p{0.24\textwidth}@{}}
\toprule
Month & Work & Milestone\\
\midrule
1 & Defects4J adapter (checkout, build, faulty methods, coverage, mutants with Major); test runner with per-class timing; the orderings; the simulator & a simulated session runs on one bug\\
2 & Repair loop with feedback and shared caches; modules M1 (Repeat Guard) and M3 (Budget Guard); flakiness checks; development set: model choice, compute estimate, all settings fixed; E1 ($2{,}000$ mutants) & \textbf{registration}; E1 done\\
3 & E2: $1{,}200$ repair sessions ($100$ bugs $\times$ $3$ sessions $\times$ $4$ orderings, each up to $20$ rounds); modules M1 and M3 replayed; E3: $150$ sessions ($50$ bugs $\times$ $3$ runs) & RQ1--RQ4 analysed\\
4 & Writing; release of the system and data & thesis draft\\
\bottomrule
\end{tabularx}
\end{table}

\section{Risks}
\label{sec:plan-risks}

\begin{table}[htbp]
\centering
\caption{Main risks and how they are handled.}
\label{tab:risks}
\footnotesize
\begin{tabularx}{\textwidth}{@{}>{\raggedright\arraybackslash}p{0.3\textwidth} >{\raggedright\arraybackslash}X@{}}
\toprule
Risk & Handling\\
\midrule
Ordering saves little & the cost breakdown (C1) stands on its own; the result is read as in Chapter~\ref{ch:outcomes}\\
\km\ only matches a simpler rule & reported as a finding, with the simpler rule recommended\\
The model service changes & version pinned, answers cached; E1 needs no model\\
Defects4J is memorised, so sessions are too short & checked on the development set; the model is chosen so that the median session needs at least three rounds\\
Audits take too long & caps per bug and a fixed order of cuts (Section~\ref{sec:exp-compute})\\
Results are noisier than assumed & more runs within budget, or the smallest detectable effect is reported\\
Flaky or order-dependent tests & detected by repeated runs and excluded\\
Four months is tight & E3 is cut first; E1 and E2 answer RQ1, RQ2 and RQ4\\
\bottomrule
\end{tabularx}
\end{table}

\section{Threats to Validity}
\label{sec:plan-threats}

\paragraph{Internal.} Costs of every ordering are computed from one audit run
in a fixed order; this assumes test results and times do not depend on order,
which is checked on a sample. Mutants are not model patches, so E1 only tests
the mechanism; every claim about the loop rests on the live runs.

\paragraph{External.} One benchmark, one language and one model; the faulty
methods are given to the model. Defects4J may be memorised by
LLMs~\citep{ramos2024memorizing}, which can make sessions shorter; this affects
every configuration alike. CI systems that run tests in parallel on many
machines are not covered. Two of the practical requirements (P2, P5) rest on
indirect evidence, and the give-up report is not tested with developers.

\paragraph{Construct.} Success means a plausible patch, not a correct one.

\section{Future Work}
\label{sec:plan-future}

\paragraph{A Python benchmark.} The first extension is to Python, using
BugsInPy~\citep{widyasari2020bugsinpy}, which is modelled on Defects4J: $493$
real bugs from $17$ Python projects, each with a faulty and a fixed version, the
tests that expose the bug, and a command-line tool to check out and run them.
Its structure lets the evaluation be repeated with few changes to the modules:
\begin{itemize}[leftmargin=*,itemsep=2pt]
\item \textbf{Module M1 (Repeat Guard)}: normalisation uses Python's own syntax
tree instead of a Java parser;
\item \textbf{Module M2 (Test Prioritiser)}: the unit of ordering becomes a
\texttt{pytest} test file, and the start times come from the timings
\texttt{pytest} reports; the Coverage baseline uses a Python coverage tool;
\item \textbf{Module M3 (Budget Guard)}: unchanged, since it reads only
failure signatures and spending.
\end{itemize}
Python has no JVM start-up or compilation step, so the share of validation
spent on running tests, and with it the room for Module M2, is expected to be
larger than on Defects4J; measuring this is itself a test of RQ1. Repository-level
tasks such as SWE-bench~\citep{jimenez2024swebench} and contest programs such
as ConDefects~\citep{wu2024condefects}, where a test is one run of the whole
program, are later steps.

\paragraph{A stronger model.} If time allows, the $20$ development bugs are
run again with the strongest tier of the same family, \texttt{gpt-6-sol}
(\$2 and \$10 per million input and output tokens), in the Default
configuration of E2 with three sessions per bug. A stronger model is expected
to need fewer rounds and to repeat itself less, which moves the cost towards
validating the accepted patch. Comparing its cost breakdown with that of the
main model shows how the room left for each module changes with the strength of
the model, a further result for RQ1, and checks that the rule for choosing the
model (Section~\ref{sec:exp-live}) does not simply pick a setting in which
the wastes exist.

\paragraph{Other directions.} Multi-file bugs and imperfect fault
localisation; a study of whether the give-up report helps developers; and using
the test record to choose which failures to show the model.

\section{Conclusion}
\label{sec:plan-conclusion}

In LLM-based repair, each round costs a model call and a test run, and part of
this cost is wasted in three ways. This proposal builds a validation layer with
one module aimed at each waste: a Repeat Guard that avoids
validating the same patch twice, a Test Prioritiser that uses \km\ to reject
wrong patches earlier, and a Budget Guard that stops unproductive sessions and
reports the reason. It also measures how the cost is distributed, which bounds
what each module can save. The layer is evaluated in a controlled
simulation on mutated patches and in live runs with a real model on Defects4J,
with every outcome's reading fixed in advance.
